\documentclass[12pt]{article}
\usepackage{amsmath,amssymb,amsfonts}
\usepackage[margin=1in]{geometry}

\begin{document}

\begin{center}
{\Large\bfseries Chapter 10, Problem 14}
\end{center}

\noindent\textbf{Problem.} Show that any $2 \times 2$ matrix which commutes with the three Pauli matrices given in Problem 10.2 must be a multiple of the unit matrix. \\[4pt]
\textit{[Hint: Show that any $2 \times 2$ matrix can be written as a linear combination of the unit matrix and the Pauli matrices.]}

\bigskip
\noindent\textbf{Solution.}

\noindent\textbf{Step 1:} Any $2 \times 2$ matrix can be written as a linear combination of $I, \sigma_x, \sigma_y, \sigma_z$.

Let $M = \begin{pmatrix} a & b \\ c & d \end{pmatrix}$ be an arbitrary $2 \times 2$ matrix. We write $M = \alpha_0 I + \alpha_1 \sigma_x + \alpha_2 \sigma_y + \alpha_3 \sigma_z$:
\[
\alpha_0 I + \alpha_1 \sigma_x + \alpha_2 \sigma_y + \alpha_3 \sigma_z =
\begin{pmatrix} \alpha_0 + \alpha_3 & \alpha_1 - i\alpha_2 \\ \alpha_1 + i\alpha_2 & \alpha_0 - \alpha_3 \end{pmatrix}.
\]

Matching entries:
\[
\alpha_0 = \frac{a+d}{2}, \quad \alpha_3 = \frac{a-d}{2}, \quad \alpha_1 = \frac{b+c}{2}, \quad \alpha_2 = \frac{c-b}{2i}.
\]

So $\{I, \sigma_x, \sigma_y, \sigma_z\}$ forms a basis for all $2\times 2$ complex matrices.

\medskip
\noindent\textbf{Step 2:} If $M$ commutes with all three Pauli matrices, then $\alpha_1 = \alpha_2 = \alpha_3 = 0$.

We use the commutation relations $[\sigma_i, \sigma_j] = 2i\epsilon_{ijk}\sigma_k$. Since $I$ commutes with everything, $[M, \sigma_j] = \alpha_1[\sigma_x, \sigma_j] + \alpha_2[\sigma_y, \sigma_j] + \alpha_3[\sigma_z, \sigma_j]$.

\textit{Commute with $\sigma_x$:}
\[
[M, \sigma_x] = \alpha_2[\sigma_y, \sigma_x] + \alpha_3[\sigma_z, \sigma_x] = \alpha_2(-2i\sigma_z) + \alpha_3(2i\sigma_y) = 0.
\]
Since $\sigma_y$ and $\sigma_z$ are linearly independent: $\alpha_3 = 0$ and $\alpha_2 = 0$.

\textit{Commute with $\sigma_y$:}
\[
[M, \sigma_y] = \alpha_1[\sigma_x, \sigma_y] = \alpha_1(2i\sigma_z) = 0 \implies \alpha_1 = 0.
\]

Therefore $M = \alpha_0 I$, i.e., $M$ is a multiple of the unit matrix. $\square$

\medskip
\noindent\textit{Remark:} This is an instance of \textbf{Schur's lemma}. The three Pauli matrices together with $I$ form a representation of the quaternion group (Problem 2). The set $\{\sigma_x, \sigma_y, \sigma_z\}$ generates (via the group in Problem 2) a two-dimensional irreducible representation. By Schur's lemma, any matrix commuting with all elements of an irreducible representation must be a scalar multiple of the identity.

\end{document}
